\documentclass[notheorems,hyperref={bookmarks=true}]{beamer}

% --- CÁC GÓI LỆNH CẦN THIẾT ---
\usepackage[framesassubsections]{beamerprosper}
\usepackage{multirow}
\usepackage[utf8]{inputenc} 
\usepackage[utf8]{vietnam} 
\usepackage{ulem}
\usepackage{pifont}
\usepackage[mathscr]{eucal}
\usepackage{amsmath, amsthm, amssymb, amsxtra, latexsym, amscd, makeidx, tikz}
\usetikzlibrary{calc, shapes.geometric, positioning}
\usepackage{graphicx}
\usepackage{array}
\usepackage{multicol}
\usepackage{multimedia}
\usepackage{xcolor}
\usepackage{booktabs} 
\usepackage[table]{xcolor}
\usetikzlibrary{shadows}

% --- CÀI ĐẶT GIAO DIỆN ---
\usetheme{CambridgeUS}
\usecolortheme{rose} 
\usefonttheme[onlylarge]{structurebold}
\usefonttheme{professionalfonts}

\setbeamercolor{title}{fg=red!80!black,bg=blue!20!white}
\setbeamertemplate{navigation symbols}{}
\setbeamertemplate{blocks}[rounded][shadow=true]

% \logo{\includegraphics[height=1cm]{logoHust.png}} % Comment nếu không có hình

\setbeamertemplate{footline}{
  \leavevmode%
  \hbox{%
    \begin{beamercolorbox}[wd=.33\paperwidth,ht=2.5ex,dp=1.5ex,center]{author in head/foot}%
      Nguyễn Tuấn Tài
    \end{beamercolorbox}%
    \begin{beamercolorbox}[wd=.34\paperwidth,ht=2.5ex,dp=1.5ex,center]{title in head/foot}%
      \footnotesize{AI Agent hỗ trợ học ngoại ngữ}
    \end{beamercolorbox}%
    \begin{beamercolorbox}[wd=.33\paperwidth,ht=2.5ex,dp=1.5ex,right]{date in head/foot}%
      \insertframenumber/\inserttotalframenumber \hspace*{2ex} 
    \end{beamercolorbox}%
  }%
  \vskip0pt%
}

\title[AI Agent dạy ngoại ngữ]{
    \Large\bfseries Xây dựng AI Agent hỗ trợ học ngoại ngữ\\
    \large Tóm tắt kết quả chính
}

\institute[\scriptsize Khoa Toán - Tin]{
    \small
    \begin{tabular}{l}
        \textbf{Giảng viên hướng dẫn: \textcolor{red!80!black}{TS. Tạ Anh Sơn}} \\
        \textbf{Sinh viên thực hiện: \textcolor{red!80!black}{Nguyễn Tuấn Tài (20216959)}} \\
    \end{tabular} \\[0.5em]
    \textbf{Khoa Toán - Tin, Đại học Bách khoa Hà Nội}
}

\date[Hà Nội, 06/2025]{\scriptsize\bfseries\color{red!80!black}Hà Nội, 06/2026}

\begin{document}

\begin{frame}
    \maketitle
\end{frame}

\footnotesize

% ==========================================
\section{Đặt vấn đề và Hướng giải quyết}
% ==========================================

\begin{frame}{Đặt vấn đề: 3 Hạn chế của phương pháp truyền thống}
\begin{block}{Vấn đề 1: Chi phí cao \& Thời gian cố định}
    \begin{itemize}
        \item Gia sư/Trung tâm: 200k-500k/giờ, lịch học cố định
        \item \textbf{→ Đồ án giải quyết:} Miễn phí, học mọi lúc mọi nơi
    \end{itemize}
\end{block}

\begin{block}{Vấn đề 2: Thiếu phản hồi cá nhân hóa}
    \begin{itemize}
        \item Ứng dụng tự học không phân tích lỗi chi tiết
        \item \textbf{→ Đồ án giải quyết:} Error Analyzer phân tích lỗi tức thì + ghi nhớ điểm yếu
    \end{itemize}
\end{block}

\begin{block}{Vấn đề 3: Nguồn luyện tập hạn chế}
    \begin{itemize}
        \item Bài tập cố định, không thích ứng độ khó
        \item \textbf{→ Đồ án giải quyết:} Exercise Generator sinh bài tập vô hạn, thích ứng trình độ
    \end{itemize}
\end{block}
\end{frame}

% ==========================================
\section{Kiến trúc hệ thống Multi-Agent}
% ==========================================

\begin{frame}{Kiến trúc Multi-Agent: 3 Agents chuyên biệt}
\begin{figure}[H]
\centering
\begin{tikzpicture}[
    node distance=1.5cm,
    agent_box/.style={
        rectangle,
        draw=gray!60,
        fill=blue!8,
        text width=4.2cm,
        align=center,
        minimum height=2.5cm,
        rounded corners=5pt,
        font=\small
    },
    arrow/.style={->, >=stealth, thick, color=gray!70}
]

% Agent 1: Error Analyzer
\node[agent_box] (ErrorAgent) at (0, 0) {
    \textbf{\color{red!80}Error Analyzer}\\[0.2cm]
    \textit{Phân tích lỗi học tập}\\[0.3cm]
    {\small \color{gray!80}
    • Nhận diện loại lỗi\\
    • Giải thích chi tiết\\
    • Đưa hướng sửa
    }
};

% Agent 2: Exercise Generator
\node[agent_box, right=of ErrorAgent, fill=green!8] (ExerciseAgent) {
    \textbf{\color{green!80}Exercise Generator}\\[0.2cm]
    \textit{Sinh bài tập tự động}\\[0.3cm]
    {\small \color{gray!80}
    • Sinh bài tập mới\\
    • Thích ứng độ khó\\
    • Luyện tập vô hạn
    }
};

% Agent 3: Writing Evaluator
\node[agent_box, right=of ExerciseAgent, fill=purple!8] (WritingAgent) {
    \textbf{\color{purple!80}Writing Evaluator}\\[0.2cm]
    \textit{Đánh giá bài viết}\\[0.3cm]
    {\small \color{gray!80}
    • Chấm 4 tiêu chí\\
    • Phản hồi chi tiết\\
    • Ghi điểm lịch sử
    }
};

% Center: AI Tutor
\node[
    rectangle,
    draw=orange!60,
    fill=orange!15,
    text width=3cm,
    align=center,
    minimum height=1.8cm,
    rounded corners=5pt,
    font=\small,
    below=2.5cm of ExerciseAgent
] (AITutor) {
    \textbf{\color{orange!80}AI Tutor}\\
    \small (Điều phối)\\[0.3cm]
    {\tiny \color{gray!80}
    Quyết định gọi\\
    agent nào khi nào
    }
};

% Arrows
\draw[arrow] (ErrorAgent.south) -- (AITutor.north west);
\draw[arrow] (ExerciseAgent.south) -- (AITutor.north);
\draw[arrow] (WritingAgent.south) -- (AITutor.north east);

\end{tikzpicture}
\caption{Kiến trúc Multi-Agent với AI Tutor điều phối 3 Agents chuyên biệt}
\end{figure}
\end{frame}

% ==========================================
\section{Kết quả triển khai}
% ==========================================

\begin{frame}{Kết quả 1: Hệ thống nội dung toàn diện}
\begin{columns}
    \begin{column}{0.45\textwidth}
        \begin{block}{190 chủ đề CEFR}
            \textbf{Mỗi chủ đề gồm:}
            \begin{enumerate}
                \item \textbf{Grammar}
                \item \textbf{Vocabulary}
                \item \textbf{Practice}
                \item \textbf{Writing}
                \item \textbf{Quiz}
            \end{enumerate}
        \end{block}
        
        \vspace{0.3cm}
        
        \begin{alertblock}{Kết quả}
            \centering
            \Large \textbf{950 bài học}\\
            \small được chuẩn hóa
        \end{alertblock}
    \end{column}
    \begin{column}{0.55\textwidth}
        \centering
        % Nếu có hình ũiuhsx.png thì dùng, không thì comment lại
        % \includegraphics[width=\textwidth]{ũiuhsx.png}
        
        % Thay bằng text mô tả
        \fbox{\begin{minipage}{0.9\textwidth}
            \centering
            \textbf{Cấu trúc lộ trình học}\\[0.3cm]
            {\footnotesize
            A1 (20 topics) → A2 (32 topics)\\
            B1 (33 topics) → B2 (36 topics)\\
            C1 (35 topics) → C2 (34 topics)\\[0.3cm]
            Tổng: \textbf{190 topics × 5 lessons}\\
            = \textbf{950 bài học}
            }
        \end{minipage}}
        
        \vspace{0.3cm}
        {\footnotesize \textit{Cung cấp ngữ cảnh phong phú cho AI Agent}}
    \end{column}
\end{columns}
\end{frame}

\begin{frame}{Kết quả 2: Phân tích lỗi thông minh}
\begin{columns}
    \begin{column}{0.45\textwidth}
        \begin{exampleblock}{Error Analyzer Agent}
            \textbf{Chức năng:}
            \begin{itemize}
                \item Phát hiện lỗi \textit{ngữ pháp/từ vựng}
                \item Giải thích bằng tiếng Việt
                \item Đưa hướng sửa cụ thể
            \end{itemize}
        \end{exampleblock}
        
        \vspace{0.3cm}
        
        \begin{alertblock}{Ưu điểm}
            \begin{itemize}
                \item Phản hồi \textbf{tức thì}
                \item Phân loại \textbf{2 cấp độ}
                \item Lưu lịch sử → \textbf{cá nhân hóa}
            \end{itemize}
        \end{alertblock}
    \end{column}
    \begin{column}{0.55\textwidth}
        \centering
        \textbf{Ví dụ phân tích lỗi:}\\[0.3cm]
        % Nếu có hình jhsavs.png hoặc error_analysis_example.png
        % \fbox{\includegraphics[width=0.9\textwidth]{jhsavs.png}}
        
        % Thay bằng ví dụ text
        \fbox{\begin{minipage}{0.9\textwidth}
            \footnotesize
            \textbf{Input:} \textit{"I \textcolor{red}{go} to school yesterday"}\\[0.2cm]
            \textbf{Lỗi:} Ngữ pháp - Thì động từ\\[0.2cm]
            \textbf{Giải thích:} "Yesterday" là dấu hiệu quá khứ, cần dùng \textbf{Past Simple}\\[0.2cm]
            \textbf{Sửa:} \textit{"I \textcolor{green}{went} to school yesterday"}\\[0.2cm]
            \textbf{Lưu vào Long-term Memory:} Điểm yếu "Past Simple"
        \end{minipage}}
    \end{column}
\end{columns}
\end{frame}

\begin{frame}{Kết quả 3: AI Tutor tương tác thông minh}
\begin{columns}
    \begin{column}{0.45\textwidth}
        \begin{exampleblock}{Chức năng}
            \textbf{1. Giải thích lý thuyết}\\
            {\footnotesize Bằng tiếng Việt, dễ hiểu}
            
            \vspace{0.2cm}
            
            \textbf{2. Tạo bài tập trong chat}\\
            {\footnotesize Khi user yêu cầu luyện tập}
            
            \vspace{0.2cm}
            
            \textbf{3. Chấm bài tức thì}\\
            {\footnotesize Feedback chi tiết + sửa lỗi}
        \end{exampleblock}
    \end{column}
    \begin{column}{0.55\textwidth}
        \centering
        \textbf{Ví dụ tương tác:}\\[0.2cm]
        \fbox{\begin{minipage}{0.9\textwidth}
            \footnotesize
            \textbf{User:} "Tôi sai Past Simple 3 lần rồi"\\[0.15cm]
            \textbf{AI Tutor:}\\
            \textcolor{blue}{• Giải thích quy tắc}\\
            \textcolor{blue}{• Cho 3 ví dụ minh họa}\\
            \textcolor{blue}{• Tạo 5 câu luyện tập}\\[0.15cm]
            \textbf{User:} "1. went 2. played..."\\[0.15cm]
            \textbf{AI Tutor:}\\
            \textcolor{green}{✅ Câu 1, 3, 5: Đúng!}\\
            \textcolor{red}{❌ Câu 2, 4: Sửa...}
        \end{minipage}}
    \end{column}
\end{columns}
\end{frame}

\begin{frame}{Kết quả 4: Đánh giá bài viết tự động}
\begin{columns}
    \begin{column}{0.5\textwidth}
        \begin{exampleblock}{Writing Evaluator Agent}
            \textbf{Chấm điểm theo 4 tiêu chí CEFR:}
            \begin{enumerate}
                \item \textbf{Ngữ pháp} (Grammar)
                \item \textbf{Từ vựng} (Vocabulary)
                \item \textbf{Nội dung} (Content)
                \item \textbf{Cấu trúc} (Structure)
            \end{enumerate}
        \end{exampleblock}
        
        \vspace{0.3cm}
        
        \begin{block}{Phản hồi chi tiết}
            \begin{itemize}
                \item Chỉ ra \textit{từng câu} cần sửa
                \item Gợi ý cách hành văn tốt hơn
                \item Lưu lịch sử → theo dõi tiến bộ
            \end{itemize}
        \end{block}
    \end{column}
    \begin{column}{0.5\textwidth}
        \centering
        \textbf{Ví dụ đánh giá bài viết:}\\[0.3cm]
        % Nếu có hình writing_evaluation_example.png thì dùng
        % \fbox{\includegraphics[width=0.9\textwidth]{writing_evaluation_example.png}}
        
        % Thay bằng bảng điểm mẫu
        \fbox{\begin{minipage}{0.9\textwidth}
            \footnotesize
            \textbf{Topic:} My Daily Routine (A1)\\[0.2cm]
            \begin{tabular}{ll}
                \textbf{Grammar:} & 7/10 \\
                \textbf{Vocabulary:} & 8/10 \\
                \textbf{Content:} & 9/10 \\
                \textbf{Structure:} & 7/10 \\
                \hline
                \textbf{Tổng:} & \textbf{7.75/10}
            \end{tabular}
            \\[0.3cm]
            \textbf{Nhận xét:} Bài viết tốt, cần chú ý thì động từ ở câu 3 và 5.
        \end{minipage}}
    \end{column}
\end{columns}
\end{frame}

\begin{frame}{Kết quả 5: Quản lý bộ nhớ học tập}
\begin{block}{Hệ thống bộ nhớ 2 tầng}
    \textbf{1. Short-term Memory (Bộ nhớ ngắn hạn)}
    \begin{itemize}
        \item Duy trì \textit{mạch hội thoại} trong phiên làm việc
        \item Ghi nhớ ngữ cảnh: chủ đề đang học, lỗi vừa mắc
    \end{itemize}
    
    \vspace{0.2cm}
    
    \textbf{2. Long-term Memory (Bộ nhớ dài hạn)}
    \begin{itemize}
        \item Lưu trữ \textit{hồ sơ lỗi sai} lâu dài (PostgreSQL)
        \item Phân tích \textit{xu hướng học tập} để cá nhân hóa
        \item Dashboard thống kê kỹ năng yếu
    \end{itemize}
\end{block}

\vspace{0.3cm}

\begin{alertblock}{Lợi ích}
    \textbf{Cá nhân hóa sâu:} AI Tutor "nhớ" điểm yếu của từng người học và ưu tiên luyện tập các kỹ năng đó
\end{alertblock}
\end{frame}

\begin{frame}{Kết quả 6: Triển khai thực tế}
\begin{table}
    \centering
    \caption{Hệ thống được triển khai hoàn chỉnh}
    \resizebox{0.95\textwidth}{!}{
        \begin{tabular}{|l|l|l|}
            \hline
            \rowcolor{blue!20} \textbf{Thành phần} & \textbf{Công nghệ} & \textbf{Trạng thái} \\ \hline
            \textbf{Backend API} & FastAPI + LangChain & ✅ Triển khai \\ \hline
            \textbf{Frontend} & Streamlit & ✅ Triển khai \\ \hline
            \textbf{Database} & PostgreSQL (Neon Tech) & ✅ Cloud \\ \hline
            \textbf{LLM Service} & Groq API (Llama 3.1 70B) & ✅ Production \\ \hline
            \textbf{Hosting} & Render.com & ✅ Live \\ \hline
        \end{tabular}
    }
\end{table}

\vspace{0.3cm}

\begin{exampleblock}{URL truy cập}
    \textbf{Website:} \texttt{https://lang-prj.onrender.com}
\end{exampleblock}

\begin{block}{Tính năng đã triển khai}
    \begin{itemize}
        \item ✅ Lộ trình học 190 chủ đề
        \item ✅ Chat với AI Tutor
        \item ✅ Phân tích lỗi tức thì
        \item ✅ Đánh giá bài viết
        \item ✅ Dashboard tiến độ
    \end{itemize}
\end{block}
\end{frame}

% ==========================================
\section{Tổng kết}
% ==========================================

\begin{frame}{Tổng kết: Các kết quả chính đạt được}
\begin{exampleblock}{Giải quyết 3 vấn đề đặt ra}
    \begin{itemize}
        \item \textbf{Vấn đề chi phí/thời gian:} ✅ Hệ thống online miễn phí, 24/7
        \item \textbf{Vấn đề cá nhân hóa:} ✅ Phân tích lỗi tức thì + bộ nhớ dài hạn
        \item \textbf{Vấn đề bài tập hạn chế:} ✅ AI Tutor tạo bài tập trong chat, hỗ trợ linh hoạt
    \end{itemize}
\end{exampleblock}

\begin{block}{6 Kết quả kỹ thuật đạt được}
    \begin{enumerate}
        \item \textbf{Kiến trúc Multi-Agent:} 3 agents chuyên biệt + AI Tutor điều phối
        \item \textbf{Mô hình tri thức:} 190 chủ đề CEFR × 5 bài học = 950 bài học
        \item \textbf{Error Analyzer:} Phân tích lỗi 2 cấp độ (Ngữ pháp/Từ vựng)
        \item \textbf{AI Tutor thông minh:} Chat tương tác, giải thích, tạo bài tập
        \item \textbf{Writing Evaluator:} Chấm điểm tự động theo 4 tiêu chí CEFR
        \item \textbf{Triển khai production:} Hệ thống live trên Cloud (Render + Neon Tech)
    \end{enumerate}
\end{block}
\end{frame}

\begin{frame}{Hạn chế và Hướng phát triển}
\begin{columns}[T]
    \begin{column}{0.5\textwidth}
        \begin{alertblock}{Hạn chế hiện tại}
            \begin{itemize}
                \item Chỉ tập trung \textbf{Đọc/Viết}
                \item Chưa có kỹ năng \textbf{Nghe/Nói}
                \item Đánh giá bài viết chưa có scoring model riêng
            \end{itemize}
        \end{alertblock}
    \end{column}
    \begin{column}{0.5\textwidth}
        \begin{exampleblock}{Hướng phát triển}
            \begin{itemize}
                \item Tích hợp \textbf{Voice AI} (Speech-to-Text, Text-to-Speech)
                \item Mở rộng \textbf{4 kỹ năng} CEFR
                \item Xây dựng \textbf{scoring model} riêng cho Writing
                \item Tối ưu chi phí LLM
            \end{itemize}
        \end{exampleblock}
    \end{column}
\end{columns}
\end{frame}

\section*{}
% ==========================================
\begin{frame}
    \begin{center}
         \Large\bfseries Cảm ơn cô và các bạn \\ đã lắng nghe!
         
         \vspace{1cm}
         

         \begin{figure}
             \centering
             \includegraphics[width=0.5\linewidth]{sloth.jpg}
         \end{figure}
         
    \end{center}
\end{frame}

\end{document}
